\subsection{Coexistence}\label{app:coexistence}
\makeatletter
\providecommand{\fixedlabel}[2]{\phantomsection\begingroup\def\@currentlabel{#2}\label{#1}\endgroup}
\makeatother

This subsection states in full and proves Proposition~\ref{prop:coexistence}, Corollary~\ref{cor:example-coexistence}
and Corollary~\ref{prop:priors}. Its first part holds against every continuous incumbent distribution with a positive
spread and \(g_H>g_L\), and along every path of such distributions; its second part applies the first to the example.

\subsubsection*{General results}

\paragraph{Setting.}
Throughout this part the classes satisfy the class ranking, and the prior \(\pi\in(0,1)\), the slot
\((\lambda,\beta)\) with \(0<\lambda<1\) and \(0<\beta<1/2\), and \(\rho\in(0,1)\), \(c_H>0\) and \(k>0\) are fixed. An
incumbent law \(F\) is continuous, with \(\DT(F)>0\) and \(g_H(F)>g_L(F)\). Against \(F\) write \(\DT\), \(\tau\) and
\(J=\DT(1-\tau)\) for the spread, the preparation threshold and the gross return of a purchase that just reaches
\(\tau\) and recruits. Write \(M=M(\pi)\), \(m=m(\pi)\),
\[
\DT^{\rm F}=\frac{k}{1-M},\qquad \DT^{\rm N}=\frac{k}{\rho\max\{\pi,1-\pi\}} .
\]
By Lemma~\ref{lem:trading}(a), no equilibrium puts weight on a wrong-signed order. Write \(s_+=\lambda\sigma_H(1)\) and
\(s_-=\lambda\sigma_L(-1)\). On the orders \((-1,0,1)\), \(f_H=(b,\;b_0+\lambda-s_+,\;b+s_+)\) and
\(f_L=(b+s_-,\;b_0+\lambda-s_-,\;b)\), so
\[
\mu_+(s_+)=\frac{\pi(b+s_+)}{b+\pi s_+},\qquad 1-\mu_+(s_+)=\frac{(1-\pi)\,b}{b+\pi s_+},\qquad
\mu_-(s_-)=\frac{\pi b}{b+(1-\pi)s_-},
\]
and \(\mu_0(s_+,s_-)\) has the odds \(\frac{\pi}{1-\pi}\cdot\frac{b_0+\lambda-s_+}{b_0+\lambda-s_-}\). On \([0,\lambda]\),
\(\mu_+\) rises strictly from \(\pi\) to \(M\), and \(\mu_-\) falls strictly from \(\pi\) to \(m\). The same formula
defines \(\mu_+(s)\) for every \(s\ge0\), rising strictly to one. For \(\pi<\tau<1\) let
\[
s^*=\frac{b\,(\tau-\pi)}{\pi\,(1-\tau)}>0 ,
\]
the rate at which a purchase reaches the threshold, \(\mu_+(s^*)=\tau\). Then \(\mu_+(s_+)<\tau\) exactly when
\(s_+<s^*\); the rate \(s^*\) is feasible, \(s^*\le\lambda\), exactly when \(\tau\le M\), and \(s^*<\lambda\) exactly when
\(\tau<M\). If \(\tau\ge M\), and in particular if \(\tau\ge1\), then \(\mu_+(s_+)<\tau\) for every \(s_+<\lambda\).

A path is a family \((F_r)_{r\in I}\) of such laws over an interval \(I\) of strengths, with \(F_{r'}\) first-order
stochastically dominating \(F_r\) for \(r'>r\) and every other primitive fixed. Along a path \(\DT(r)\) is
nondecreasing, \(g_H(r)\), \(g_L(r)\) and \(\gamma(r)=g_H(r)-g_L(r)\) are nonincreasing, and \(\DT(r)+\gamma(r)\) is
constant (Lemma~\ref{lem:twonumbers}). A set of strengths is an up-set if it contains every strength above each of
its members, and a down-set if it contains every strength below each of its members.

The posteriors follow from (E1), because the low class never buys and class \(H\) never sells. The posterior
\(\mu_+(s)\) is continuous in \(s\ge0\), with derivative of the sign of \(\pi b(1-\pi)>0\), value \(\pi\) at \(s=0\),
value \(M\) at \(s=\lambda\) and limit one, and the equation \(\mu_+(s)=\tau\) is the linear equation
\(\pi(b+s)=\tau(b+\pi s)\), whose unique root is \(s^*\). The statements on \(s^*\) follow. Entry after each order
follows from Lemma~\ref{lem:prices}(d) at its posterior.

\begin{lemma}[The high class]\label{lem:a4-high}
Let \(\pi<\tau\). In a candidate with no wrong-signed order, the pair \((s_+,e(1))\) satisfies (E1) to (E3) at the
order \(1\), and (E5) for class \(H\), exactly when it is one of these five pairs.
\begin{enumerate}[leftmargin=2.5em]
\item[(N)] \(s_+=0\) and \(e(1)=\rho\), with \(\rho\DT(1-\pi)\le k\).
\item[(Q)] \(0<s_+<\lambda\), \(\mu_+(s_+)<\tau\), \(e(1)=\rho\) and \(\rho\DT[1-\mu_+(s_+)]=k\). This pair needs an
indifferent investor who may mix.
\item[(R)] \(0<s_+<\lambda\), \(\mu_+(s_+)>\tau\), \(e(1)=1\) and \(\DT[1-\mu_+(s_+)]=k\). This pair needs an
indifferent investor who may mix.
\item[(T)] \(0<s_+<\lambda\), \(\mu_+(s_+)=\tau\), that is \(s_+=s^*\), and \(e(1)=e^*:=k/J\). When an indifferent
high-cost challenger prepares, this requires \(e^*=1\), that is \(J=k\). When an indifferent high-cost challenger
may prepare with any probability, it requires \(\rho\le e^*\le1\), that is \(\rho J\le k\le J\). This pair needs an
indifferent investor who may mix.
\item[(F)] \(s_+=\lambda\). Here \(e(1)=1\) if \(M>\tau\) and \(e(1)=\rho\) if \(M<\tau\). If \(M=\tau\), then \(e(1)=1\)
when an indifferent high-cost challenger prepares, and \(e(1)\in[\rho,1]\) when it may prepare with any probability.
The return satisfies \(e(1)\DT(1-M)\ge k\) when an indifferent investor may mix, and \(e(1)\DT(1-M)>k\) when an
indifferent investor abstains.
\end{enumerate}
Each equation in (Q) and (R) has at most one solution. If \(\tau\ge M\), and in particular if \(\tau\ge1\), the pairs
(R) and (T) do not occur, and (F) has \(e(1)=\rho\) when \(\tau>M\).
\end{lemma}

\begin{proof}
By Lemma~\ref{lem:trading}(b), (E5) for class \(H\) requires \(K_H\le k\) at \(s_+=0\). At \(s_+=\lambda\) it requires
\(K_H\ge k\) when an indifferent investor may mix, and \(K_H>k\) when an indifferent investor abstains, because that
rule selects abstention at a tie. At \(0<s_+<\lambda\) it requires \(K_H=k\); an interior rate is a mixture at
indifference, so it needs an indifferent investor who may mix. Here \(K_H=e(1)\DT[1-\mu_+(s_+)]\), and \(e(1)\)
follows from Lemma~\ref{lem:prices}(d) at the posterior \(\mu_+(s_+)\).
\begin{itemize}[leftmargin=2em]
\item At \(s_+=0\), \(\mu_+(0)=\pi<\tau\), so \(e(1)=\rho\). This is (N).
\item At \(0<s_+<\lambda\) with \(\mu_+(s_+)<\tau\): \(e(1)=\rho\), and \(K_H=k\) is the equation of (Q).
\item At \(0<s_+<\lambda\) with \(\mu_+(s_+)>\tau\): \(e(1)=1\), and \(K_H=k\) is the equation of (R).
\item At \(0<s_+<\lambda\) with \(\mu_+(s_+)=\tau\): then \(\pi<\tau<M<1\) and \(s_+=s^*\), so \(J>0\), and
\(K_H=e(1)\DT(1-\tau)=e(1)J=k\) gives \(e(1)=e^*\). When an indifferent high-cost challenger prepares,
\(e(1)=1\); when it may prepare with any probability, \(e(1)\) can take any value in \([\rho,1]\). This is (T).
\item At \(s_+=\lambda\), \(\mu_+(\lambda)=M\). Lemma~\ref{lem:prices}(d) gives \(e(1)\), and the condition on
\(K_H\) is that of (F).
\end{itemize}
The function \(1-\mu_+\) falls strictly and \(\DT>0\), so each equation has at most one root. If \(\tau\ge M\), then
\(\mu_+(s_+)<M\le\tau\) for every \(s_+<\lambda\), so (R) and (T) do not occur, and \(M<\tau\) gives \(e(1)=\rho\) in (F).
\end{proof}

\begin{lemma}[The low class]\label{lem:a4-low}
Let \(\pi<\tau\). Then \(\mu_-(s_-)\le\pi<\tau\) for every \(s_-\), so \(e(-1)=\rho\) in every candidate. In a
candidate with no wrong-signed order, the rate \(s_-\) satisfies (E1) to (E3) at the order \(-1\), and (E5) for class
\(L\), exactly when it is one of these:
\begin{enumerate}[leftmargin=2.5em]
\item[(N)] \(s_-=0\), with \(\rho\DT\pi\le k\);
\item[(Q)] \(0<s_-<\lambda\), with \(\rho\DT\mu_-(s_-)=k\); this rate needs an indifferent investor who may mix;
\item[(F)] \(s_-=\lambda\), with \(\rho\DT m\ge k\) when an indifferent investor may mix, and \(\rho\DT m>k\) when an
indifferent investor abstains.
\end{enumerate}
When an indifferent investor may mix, exactly one rate qualifies:
\[
\hat s_-=\begin{cases}0, & \rho\DT\pi\le k,\\[2pt]
\dfrac{b\,(\pi\rho\DT/k-1)}{1-\pi}, & \rho\DT m<k<\rho\DT\pi,\\[6pt]
\lambda, & \rho\DT m\ge k.\end{cases}
\]
When an indifferent investor abstains, no rate qualifies if \(\rho\DT m\le k<\rho\DT\pi\).
\end{lemma}

\begin{proof}
Since \(\mu_-\) is decreasing, \(\mu_-(s_-)\le\mu_-(0)=\pi<\tau\), so Lemma~\ref{lem:prices}(d) gives \(e(-1)=\rho\).
Then \(K_L=\rho\DT\mu_-(s_-)\) is continuous and strictly decreasing in \(s_-\), from \(\rho\DT\pi\) at \(s_-=0\) to
\(\rho\DT m\) at \(s_-=\lambda\). By Lemma~\ref{lem:trading}(b), (E5) for class \(L\) requires \(K_L\le k\) at
\(s_-=0\); \(K_L\ge k\) at \(s_-=\lambda\) when an indifferent investor may mix, and \(K_L>k\) when it abstains; and
\(K_L=k\) at \(0<s_-<\lambda\), which needs an indifferent investor who may mix. If \(\rho\DT\pi\le k\), then
\(K_L<k\) for every \(s_->0\), so only \(s_-=0\) qualifies. If \(\rho\DT m\ge k\), then \(K_L>k\) for every
\(s_-<\lambda\), so only \(s_-=\lambda\) qualifies. Otherwise \(\rho\DT m<k<\rho\DT\pi\), and the intermediate value
theorem and strict monotonicity give exactly one interior root. The equation \(\pi b/[b+(1-\pi)s_-]=k/(\rho\DT)\)
gives the formula. When an indifferent investor abstains and \(\rho\DT m\le k<\rho\DT\pi\), the rate \(s_-=0\) fails
because \(K_L>k\) there; the rate \(s_-=\lambda\) fails because \(K_L\le k\) there; and an interior rate needs an
indifferent investor who may mix.
\end{proof}

\begin{lemma}[Product structure]\label{lem:a4-product}
Let \(\pi<\tau\), and take rates with \(s_+\ge s_-\). Then \(\mu_0(s_+,s_-)\le\pi<\tau\). Under each pair of tie rules,
the tuple with the rates \((s_+,s_-)\) is an equilibrium if and only if these four conditions hold:
\begin{enumerate}[label=(\roman*),leftmargin=2.5em]
\item \((s_+,e(1))\) is a pair of Lemma~\ref{lem:a4-high};
\item \(s_-\) is a rate of Lemma~\ref{lem:a4-low};
\item \(e(0)=e(-1)=\rho\);
\item the prices are \(P(x)=t_0+e(x)\,w(\mu(x))\), and \(\bar\mu(P(x))=\mu(x)\).
\end{enumerate}
\end{lemma}

\begin{proof}
If \(s_+\ge s_-\), then \(b_0+\lambda-s_+\le b_0+\lambda-s_-\), so the odds of \(\mu_0\) are at most \(\pi/(1-\pi)\). Hence
\(\mu_0\le\pi<\tau\).

\emph{Necessity.} In any equilibrium, Lemma~\ref{lem:prices}(d) gives \(e(0)=e(-1)=\rho\) at the posteriors \(\mu_0\)
and \(\mu_-(s_-)\), both below \(\tau\), and Lemma~\ref{lem:prices}(b)--(c) give condition (iv).
Lemma~\ref{lem:trading}(d) shows that (E5) for class \(H\) involves only \((s_+,\mu(1),e(1))\), and (E5) for class
\(L\) only \((s_-,\mu(-1),e(-1))\). So (E5) gives (i) and (ii).

\emph{Sufficiency.} Take rates and entry that satisfy (i) to (iv). Each condition of
Definition~\ref{def:equilibrium} holds at this tuple. (E1): the posteriors are computed from \(f_H\) and \(f_L\).
(E2): the stated entry is \(\rho\) at the posteriors \(\mu_-(s_-)\) and \(\mu_0\), both at most \(\pi<\tau\), and it is
\(e(1)\ge\rho\) at \(\mu_+(s_+)\ge\pi\), with \(e(1)=1\) when \(\mu_+(s_+)>\tau\). So entry lies in \([\rho,1]\) and is
nondecreasing in the posterior. Two orders with a common posterior include the order \(0\) or \(-1\), so the common
posterior is at most \(\pi\); if the order \(1\) is one of them, then \(\mu_+(s_+)\le\pi\) forces \(s_+=0\), the pair
(N), with \(e(1)=\rho\). Orders with a common posterior therefore have a common entry. By the paragraph on
constructed tuples in Appendix~\ref{app:beliefs}, orders share a price exactly when they share the posterior and
the entry, the preparation rule \(\zeta(P(x))=[e(x)-\rho]/(1-\rho)\) is a function of the price, and the belief at
each price is the common posterior of its orders, which is (E2). (E3): at each price the belief is the posterior of
its orders, and the stated entry is the one that Lemma~\ref{lem:prices}(d) allows at that posterior under the
challenger rule in use; condition (i) states it at the order \(1\), and the posteriors \(\mu_-(s_-)\) and \(\mu_0\)
lie below \(\tau\). (E4): \eqref{eq:a2-terminal} with the stated entry. (E5): by Lemma~\ref{lem:trading}(d), the
payoffs of class \(H\) involve only \((s_+,\mu(1),e(1))\), and those of class \(L\) only \((s_-,\mu(-1),e(-1))\);
conditions (i) and (ii) state optimality for each, and Lemma~\ref{lem:trading}(a) rules out the wrong-signed
orders. (E6): Appendix~\ref{app:payoffs}.
\end{proof}

% STRUCTURE: full statement of Proposition 3 <- Section 4.3
\begin{named}{Proposition~\ref{prop:coexistence} (Quiet and informative markets, full statement)}\fixedlabel{prop:coexistence-app}{3}
Let the classes satisfy the class ranking, and fix the prior, the slot, \(\rho\), \(c_H\) and \(k\) as above. Parts
(a) to (c), except the second sentence of (c)(iv), take the default tie rules: an indifferent investor may mix over
its best orders, and an indifferent high-cost challenger prepares. That sentence and part (d) state their own.
\begin{enumerate}[label=(\alph*),leftmargin=2em]
\item \emph{The region.} A quiet market and an informative market coexist against \(F\) if and only if
\begin{equation}\label{eq:region}
\pi<\tau\le M\qquad\text{and}\qquad \frac{k}{1-\tau}\le\DT\le\DT^{\rm N} .
\end{equation}
Against such an \(F\) the equilibria are exactly the quiet market and one informative market. In the informative
market the low class abstains, class \(H\) buys at the rate
\(s_+=\min\bigl\{\lambda,\;(b/\pi)\,[(1-\pi)\DT/k-1]\bigr\}\), a purchase recruits a high-cost challenger, and entry
after the other two orders is \(\rho\). It is a full purchase if \(\DT\ge\DT^{\rm F}\), the threshold pair with entry
one after a purchase if \(\DT<\DT^{\rm F}\) and \(J=k\), and a mixed purchase that recruits if \(\DT<\DT^{\rm F}\) and
\(J>k\). Its entry is
\begin{equation}\label{eq:entry-inf}
\Ent=\rho+(1-\rho)\min\Bigl\{\frac{(1-\pi)\,b\,\DT}{k},\;b+\pi\lambda\Bigr\},
\end{equation}
which depends on \(F\) only through the spread, rises strictly with the spread on \(\DT<\DT^{\rm F}\) and is constant
on \(\DT\ge\DT^{\rm F}\). The set of pairs \((\DT,\tau)\) that satisfy \eqref{eq:region} is convex, and each of its
points has \(\tau\le1-\rho\max\{\pi,1-\pi\}\). Coexistence therefore requires \(\pi<1/(1+\rho)\).
\item \emph{The sure-purchase stretch.} Along a path, the strengths at which a quiet market coexists with a full
purchase, those with \(\pi<\tau(r)\le M\) and \(\DT^{\rm F}\le\DT(r)\le\DT^{\rm N}\), form an interval. Every strength at
which a quiet market coexists with a mixed purchase that recruits, or with the threshold pair, lies below every
strength of this interval. If a quiet market and an informative market coexist at \(r'\) and at \(r''\) but not at
\(r\), with \(r'<r<r''\), then \(\DT(r)<\DT^{\rm F}\): a gap lies at spreads too narrow for a sure purchase to pay.
\item \emph{The balance condition.} Along a path, let the product \(\DT(r)\,[g_H(r)-c_H]\) be nondecreasing in \(r\).
\begin{enumerate}[label=(\roman*),leftmargin=2em]
\item The strengths with \(J(r)\ge k\) form an up-set, on which \(J\) is nondecreasing and rises strictly wherever
\(\DT\) rises strictly. The strengths at which a quiet market and an informative market coexist form an interval,
at every prior and every slot.
\item Let \(I\) be compact, let \(\DT\) and \(\tau\) be continuous in \(r\), and let \(\tau>\pi\) at the weakest
incumbent of the path. Then this interval is closed. When it is not empty, it starts at the weakest strength with
\(J(r)\ge k\) and ends at the smaller of the strongest strength with \(\tau(r)\le M\) and the strongest strength with
\(\DT(r)\le\DT^{\rm N}\).
\item The product \(\DT(r)\,[g_H(r)-c_H]\) is nondecreasing, and \(J\) rises strictly, when the incumbent is uniform,
\(F_r=U[0,r]\) for \(r\) in a compact interval \(I\) with \(p<\min I\), every low-class value lies in
\([p,\min I]\), every high-class value is at least \(\max I\), and \(\E_H\theta\ge\max I+c_H\).
\item If, in addition, \(\DT(r)\,[g_H(r)-c_H]^2\) is nondecreasing, then \(\DT(r)\,[1-\tau(r)]^2\) is nondecreasing.
When an indifferent investor may mix and an indifferent high-cost challenger may prepare with any probability, the
threshold pair is a pair of Lemma~\ref{lem:a4-high} exactly at the strengths with \(\pi<\tau(r)<M\) and
\(\rho J(r)\le k\le J(r)\); on these strengths its entry, with entry \(\rho\) after a sale and after the zero order,
\[
\Ent_{\rm T}(r)=\rho+(1-\pi)\,b\,\frac{k-\rho J(r)}{\DT(r)\,[1-\tau(r)]^2},
\]
is nonincreasing. Where \(\tau(r)=M\), the threshold rate is \(s^*=\lambda\), and a purchase at that rate is a full
purchase, the pair (F) of Lemma~\ref{lem:a4-high}.
\end{enumerate}
\item \emph{Existence at the prior \(1/2\).} Let the prior be \(1/2\) and \(\tau>1/2\), including \(\tau\ge1\), and let an
indifferent investor mix over its best orders. Under either challenger rule an equilibrium exists against \(F\). If
\(\rho\DT\le2k\), the quiet market is one. If \(\rho\DT>2k\), the low-class rate of Lemma~\ref{lem:a4-low} is
\(\hat s_-=\min\{\lambda,\,b\,(\rho\DT/k-2)\}>0\). If \(\mu_+(\hat s_-)<\tau\), both classes trade at the rate
\(\hat s_-\) and entry is \(\rho\) after every order. If \(\mu_+(\hat s_-)\ge\tau\), the low class sells at the rate
\(\hat s_-\), class \(H\) buys at the rate \(\min\{\lambda,\,b\,(\DT/k-2)\}\ge\hat s_-\), a purchase recruits a
high-cost challenger, and entry after a sale and after the zero order is \(\rho\).
\end{enumerate}
\end{named}

A quiet market is an equilibrium with \(\sigma_H=\sigma_L=\delta_0\), and an informative market is an equilibrium
with informed trade.

\begin{proof}[Proof of Proposition~\ref{prop:coexistence}(a)]
Fix \(F\) and the default tie rules. The proof has four steps.

\emph{Step 1: a prior at or above the threshold.} Let \(\pi\ge\tau\). In a quiet market every posterior is
\(\pi\ge\tau\), so Lemma~\ref{lem:prices}(d) gives entry one after every order, and Lemma~\ref{lem:trading} gives
\(K_H=\DT(1-\pi)\) and \(K_L=\DT\pi\). A quiet market therefore needs \(\DT\max\{\pi,1-\pi\}\le k\). Under this
inequality no candidate with informed trade is an equilibrium. By Lemma~\ref{lem:trading}(a) such a candidate puts no
weight on a wrong-signed order. If \(s_+>0\), then \(\mu(1)=\mu_+(s_+)>\pi\ge\tau\), so \(e(1)=1\) and
\(K_H=\DT[1-\mu_+(s_+)]<\DT(1-\pi)\le k\): class \(H\) strictly prefers abstention. If \(s_->0\), then
\(\mu(-1)=\mu_-(s_-)<\pi\) and \(e(-1)\le1\), so \(K_L<\DT\pi\le k\): class \(L\) strictly prefers abstention. So
coexistence needs \(\pi<\tau\).

\emph{Step 2: the quiet market when \(\pi<\tau\).} Let \(\pi<\tau\), so Lemmas~\ref{lem:a4-high} to~\ref{lem:a4-product}
apply. A quiet market has \(s_+=s_-=0\), so by Lemma~\ref{lem:a4-product} it is an equilibrium exactly when
\((0,e(1))\) is the pair (N) of Lemma~\ref{lem:a4-high} and \(0\) is the rate (N) of Lemma~\ref{lem:a4-low}, that is
when \(\rho\DT(1-\pi)\le k\) and \(\rho\DT\pi\le k\), or \(\DT\le\DT^{\rm N}\). Its entry is then \(\rho\) after every
order and its prices follow from condition (iv), so there is exactly one quiet market.

\emph{Step 3: informative markets when \(\pi<\tau\) and \(\DT\le\DT^{\rm N}\).} Since \(\rho\DT\pi\le k\),
Lemma~\ref{lem:a4-low} gives \(s_-=0\) in every equilibrium, so \(s_+\ge s_-\), and by Lemma~\ref{lem:a4-product} the
equilibria are the pairs of Lemma~\ref{lem:a4-high} with \(s_-=0\) and entry \(\rho\) after the orders \(0\) and
\(-1\). An informative market has \(s_+>0\). The pair (Q) is impossible, because
\(\rho\DT[1-\mu_+(s_+)]<\rho\DT(1-\pi)\le k\) for \(s_+>0\), and so is (F) with \(M<\tau\), for the same reason. The
remaining pairs, (R), (T) and (F) with \(M\ge\tau\), need \(\mu_+(s_+)\ge\tau\) for some \(s_+\le\lambda\), so
\(\tau\le M\); since \(e(1)\le1\), they also need \(k\le e(1)\DT[1-\mu_+(s_+)]\le\DT(1-\tau)=J\), that is
\(\DT\ge k/(1-\tau)\). With Step 1, coexistence implies \eqref{eq:region}.

Conversely, let \(F\) satisfy \eqref{eq:region}. Then \(\tau\le M<1\) and \(J\ge k\), and Step 2 gives the quiet
market. If \(\DT\ge\DT^{\rm F}\), that is \(\DT(1-M)\ge k\), the pair (F) holds with \(e(1)=1\), which an indifferent
high-cost challenger chooses when \(M=\tau\). It is the only informative pair: at every \(s_+<\lambda\) with
\(\mu_+(s_+)\ge\tau\) the return \(\DT[1-\mu_+(s_+)]\) exceeds \(\DT(1-M)\ge k\), so neither (R) nor (T) holds. If
\(\DT<\DT^{\rm F}\), then \(\DT(1-M)<k\le J=\DT(1-\tau)\), so \(\tau<M\) and (F) fails. The return
\(s\mapsto\DT[1-\mu_+(s)]\) is continuous and strictly decreasing on \([s^*,\lambda]\), from \(J\ge k\) to
\(\DT(1-M)<k\), so it equals \(k\) at exactly one \(s_+\in[s^*,\lambda)\). At \(s_+=s^*\), which happens exactly when
\(J=k\), this is (T) with \(e^*=1\); at \(s_+>s^*\) it is (R); and no other informative pair exists. In every case the
low class abstains, entry after the orders \(0\) and \(-1\) is \(\rho\), and Lemma~\ref{lem:a4-product} completes the
tuple. The equation \(\DT(1-\pi)b/(b+\pi s)=k\) has the root \(s=(b/\pi)[(1-\pi)\DT/k-1]\), which is at least
\(\lambda\) exactly when \(\DT(1-M)\ge k\), because \(1-M=(1-\pi)b/(b+\pi\lambda)\). This gives the stated rate.

\emph{Step 4: entry and the shape of the region.} Entry is \(\Ent=\sum_x\Pr(X=x)\,e(x)\). With \(e(0)=e(-1)=\rho\),
\(e(1)=1\) and \(\Pr(X=1)=\pi(b+s_+)+(1-\pi)b=b+\pi s_+\), entry is \(\rho+(1-\rho)(b+\pi s_+)\), and
\(b+\pi s_+=\min\{(1-\pi)b\DT/k,\,b+\pi\lambda\}\) gives \eqref{eq:entry-inf}. The prior, \(b\), \(\lambda\) and \(k\)
are fixed, so entry depends on \(F\) only through \(\DT\). The two terms of the minimum are equal exactly at
\(\DT=\DT^{\rm F}\), so entry rises strictly with \(\DT\) below \(\DT^{\rm F}\) and is constant above. With
\(\tau\le M<1\), the inequality \(k/(1-\tau)\le\DT\) is \(\tau\le1-k/\DT\). The set of pairs \((\DT,\tau)\) that satisfy
\eqref{eq:region} is therefore the intersection of the hypograph of the concave function \(\DT\mapsto1-k/\DT\) on
\(\DT>0\) with three half-planes, and it is convex. Its two spread bounds give \(1-\tau\ge\rho\max\{\pi,1-\pi\}\). So
\(\pi<\tau\le1-\rho\max\{\pi,1-\pi\}\le1-\rho\pi\), which requires \(\pi<1/(1+\rho)\).
\end{proof}

\begin{proof}[Proof of Proposition~\ref{prop:coexistence}(b)]
Along a path, \(B_r(\mu)=g_L(r)+\mu\gamma(r)=\mu g_H(r)+(1-\mu)g_L(r)\) is nonincreasing in \(r\) at every
\(\mu\in[0,1]\), because \(g_H\) and \(g_L\) are, and \(\DT\) is nondecreasing (Lemma~\ref{lem:twonumbers}). Since
\(\gamma>0\), the condition \(\pi<\tau(r)\) is \(B_r(\pi)<c_H\), an up-set, and \(\tau(r)\le M\) is \(B_r(M)\ge c_H\), a
down-set; \(\DT(r)\ge\DT^{\rm F}\) is an up-set, and \(\DT(r)\le\DT^{\rm N}\) a down-set. An intersection of up-sets and
down-sets is an interval: if \(r'<r<r''\) and \(r'\) and \(r''\) belong to it, then \(r\) belongs to each up-set
because \(r'\) does, and to each down-set because \(r''\) does.

By part (a), a quiet market coexists with a full purchase exactly where \eqref{eq:region} holds and
\(\DT(r)\ge\DT^{\rm F}\). There \(J=\DT(1-\tau)\ge\DT(1-M)\ge k\), so the condition \(k/(1-\tau)\le\DT\) holds, and the
set is \(\{r:\pi<\tau(r)\le M,\ \DT^{\rm F}\le\DT(r)\le\DT^{\rm N}\}\), an interval. By part (a), a strength \(r\) at which
a quiet market coexists with a mixed purchase that recruits or with the threshold pair has
\(\DT(r)<\DT^{\rm F}\le\DT(r')\) at every strength \(r'\) of this interval, so \(r<r'\), because \(\DT\) is
nondecreasing. Finally, \(\{r:\pi<\tau(r)\le M,\ \DT(r)\le\DT^{\rm N}\}\) is an interval that contains every strength
with coexistence. If coexistence holds at \(r'\) and \(r''\) but not at \(r\in(r',r'')\), then \(r\) lies in this
interval, and \(\DT(r)\ge\DT^{\rm F}\) would put \(r\) in the set of full purchases; so \(\DT(r)<\DT^{\rm F}\).
\end{proof}

\begin{proof}[Proof of Proposition~\ref{prop:coexistence}(c)(i) and~(c)(ii)]
Since \(1-\tau=(g_H-c_H)/\gamma\), \(J=\DT\,(g_H-c_H)/\gamma\) at every strength, with \(\gamma>0\) nonincreasing and
\(\gamma=\bar\Gamma-\DT\) for the fixed class mean gap \(\bar\Gamma=\E_H\theta-\E_L\theta\) (Lemma~\ref{lem:twonumbers}).
Let \(J(r)\ge k\) and \(r'>r\). Then \(\DT(r)[g_H(r)-c_H]=\gamma(r)J(r)>0\), and the balance condition gives
\[
J(r')=\frac{\DT(r')[g_H(r')-c_H]}{\gamma(r')}\ge\frac{\DT(r)[g_H(r)-c_H]}{\gamma(r')}
\ge\frac{\DT(r)[g_H(r)-c_H]}{\gamma(r)}=J(r)\ge k .
\]
So \(\{J\ge k\}\) is an up-set on which \(J\) is nondecreasing. If \(\DT(r')>\DT(r)\), then \(\gamma(r')<\gamma(r)\) and the
second inequality is strict. By part (a), and because \(\tau\le M<1\) at every strength with coexistence, the
strengths with coexistence are \(\{\pi<\tau\}\cap\{J\ge k\}\cap\{\tau\le M\}\cap\{\DT\le\DT^{\rm N}\}\): two up-sets
and two down-sets, by the proof of part (b). They form an interval, and the argument holds at every prior and every
slot.

For (c)(ii), \(\tau>\pi\) at the weakest strength and \(\{\pi<\tau\}\) is an up-set, so it is all of \(I\). Continuity
of \(\DT\) and \(\tau\) makes \(\{J\ge k\}\), \(\{\tau\le M\}\) and \(\{\DT\le\DT^{\rm N}\}\) closed subsets of the compact
interval \(I\). A closed up-set of \(I\) is empty or runs from its least element to \(\max I\), and a closed down-set
is empty or runs from \(\min I\) to its largest element. The set with coexistence is therefore empty or the closed
interval from the weakest strength with \(J\ge k\) to the smaller of the strongest strength with \(\tau\le M\) and
the strongest strength with \(\DT\le\DT^{\rm N}\).
\end{proof}

\begin{proof}[Proof of Proposition~\ref{prop:coexistence}(c)(iii)]
The laws \(U[0,r]\), \(r\in I\), are continuous and ordered by first-order stochastic dominance. Write
\(\bar\Gamma=\E_H\theta-\E_L\theta\) for the class mean gap. Every \(r\in I\) has \(r\ge\min I>p\), so
\((r^2-p^2)/(2r^2)>0\) on \(I\). The proof has four steps.

\emph{Step 1: the auction terms.} Fix \(r\in I\) and let \(R\) be uniform on \([0,r]\). A high-class value
\(\theta\ge\max I\ge r\) wins against every incumbent value and pays \(\max\{p,R\}\), whose mean is
\(r/2+p^2/(2r)\). A low-class value \(\theta\in[p,\min I]\subseteq[p,r]\) earns
\(\E(\theta-\max\{p,R\})_+=(\theta^2-p^2)/(2r)\), and the target receives \(\theta\) minus this rent. Averaging over
each class gives
\[
g_H=\E_H\theta-\frac r2-\frac{p^2}{2r},\qquad g_L=\frac{\E_L\theta^2-p^2}{2r},\qquad
\DT=\frac r2+\frac{p^2}{2r}-\E_L\theta+\frac{\E_L\theta^2-p^2}{2r}=\frac{\E_L(r-\theta)^2}{2r},
\]
so \(\gamma=g_H-g_L=\E_H\theta-r/2-\E_L\theta^2/(2r)\) and \(\DT+\gamma=\bar\Gamma\) on \(I\). These functions are smooth
in \(r>0\), and
\[
\DT'=\frac{r^2-\E_L\theta^2}{2r^2}\ge0,\qquad g_H'=-\frac{r^2-p^2}{2r^2},\qquad \gamma'=-\DT',
\]
where \(\DT'\ge0\) because \(\E_L\theta^2\le(\min I)^2\le r^2\). Also \(\DT(r)>0\) for every \(r>\min I\), because every
low-class value is at most \(\min I\).

\emph{Step 2: three inequalities.} First, \(\E_H\theta-c_H\ge\max I\ge r\), so
\[
g_H-c_H=\E_H\theta-c_H-\frac r2-\frac{p^2}{2r}\ge\frac r2-\frac{p^2}{2r}=r\,\frac{r^2-p^2}{2r^2}>0 .
\]
Second, every low-class value lies in \([0,r]\), so
\[
r\DT'-\DT=\E_L\theta-\frac{\E_L\theta^2}{r}=\E_L\Bigl[\theta\Bigl(1-\frac\theta r\Bigr)\Bigr]\ge0 .
\]
Third, \(\E_H\theta\ge\max I+c_H\ge r+c_H\) and \(\E_L\theta^2\le r^2\), so \(\gamma\ge(r+c_H)-r/2-r^2/(2r)=c_H>0\).
Hence \(g_H>g_L\) and \(\bar\Gamma=\DT+\gamma>0\).

\emph{Step 3: the product.} By Steps 1 and 2,
\[
\bigl(\DT\,(g_H-c_H)\bigr)'=\DT'\,(g_H-c_H)-\frac{r^2-p^2}{2r^2}\,\DT\ge\frac{r^2-p^2}{2r^2}\,(r\DT'-\DT)\ge0 ,
\]
so \(\DT(r)\,[g_H(r)-c_H]\) is nondecreasing on \(I\).

\emph{Step 4: the return \(J\).} Since \(J=\DT\,(g_H-c_H)/\gamma\) and \(\gamma'=-\DT'\),
\begin{align*}
J'&=\frac{\bigl(\DT\,(g_H-c_H)\bigr)'\gamma+\DT\,(g_H-c_H)\,\DT'}{\gamma^2}
=\frac{\bar\Gamma\,(g_H-c_H)\,\DT'-\frac{r^2-p^2}{2r^2}\,\DT\,\gamma}{\gamma^2}\\
&\ge\frac{r^2-p^2}{2r^2}\cdot\frac{\bar\Gamma\,r\DT'-\DT\,\gamma}{\gamma^2}
\ge\frac{r^2-p^2}{2r^2}\cdot\frac{\DT\,(\bar\Gamma-\gamma)}{\gamma^2}=\frac{(r^2-p^2)\,\DT^2}{2r^2\gamma^2}.
\end{align*}
The first inequality uses \(g_H-c_H\ge r\,(r^2-p^2)/(2r^2)\) and \(\DT'\ge0\); the second uses \(r\DT'\ge\DT\) and
\(\bar\Gamma>0\). The bound is nonnegative on \(I\) and positive wherever \(\DT>0\), which includes every \(r>\min I\) by
Step 1, so \(J\) is strictly increasing on \(I\).
\end{proof}

For point-mass classes at \(h\) and \(\ell\), the hypotheses of part (c)(iii) read \(\ell\in[p,\min I]\),
\(p<\min I\) and \(\max I+c_H\le h\).

\begin{proof}[Proof of Proposition~\ref{prop:coexistence}(c)(iv)]
Since \(1-\tau=(g_H-c_H)/\gamma\), \(\DT(1-\tau)^2=\DT(g_H-c_H)^2/\gamma^2\). The numerator is nonnegative and
nondecreasing by hypothesis, and \(1/\gamma^2\) is positive and nondecreasing, because \(\gamma>0\) is nonincreasing;
their product is nondecreasing.

By the pair (T) of Lemma~\ref{lem:a4-high}, when an indifferent investor may mix and an indifferent high-cost
challenger may prepare with any probability, the threshold pair needs \(0<s^*<\lambda\), that is \(\pi<\tau<M\), and
\(\rho\le k/J\le1\), that is \(\rho J\le k\le J\). With entry \(\rho\) after a sale and after the zero order, its entry is
\(\rho+(k/J-\rho)(b+\pi s^*)\), and \(b+\pi s^*=(1-\pi)b/(1-\tau)\), which gives the display. Let \(r<r'\) be two
strengths of this set. Both have \(J\ge k\), so part (c)(i) gives \(J(r')\ge J(r)\) and
\(0\le k-\rho J(r')\le k-\rho J(r)\), and the first sentence gives \(\DT(r')[1-\tau(r')]^2\ge\DT(r)[1-\tau(r)]^2>0\).
Hence \(\Ent_{\rm T}(r')\le\Ent_{\rm T}(r)\). Finally, \(s^*=\lambda\) exactly when \(\tau=M\), and the purchase at the
rate \(\lambda\) is the pair (F).
\end{proof}

\begin{proof}[Proof of Proposition~\ref{prop:coexistence}(d)]
Fix \(F\) and the slot, and let the prior be \(1/2\). Since \(\pi<\tau\), Lemmas~\ref{lem:a4-high} to~\ref{lem:a4-product}
apply against \(F\). The proof has five steps.

\emph{Step 1: the frame.} At the prior \(1/2\),
\[
1-\mu_+(s)=\frac{b}{2b+s}=\mu_-(s)\qquad\text{for every }s\in[0,\lambda],
\]
and in particular \(1-M=m=b/(2b+\lambda)\). Each tuple below has no wrong-signed order, rates with \(s_+\ge s_-\),
entry \(\rho\) after a sale and after the zero order, a stated entry \(e(1)\) after a purchase, the prices
\(P(x)=t_0+e(x)\,w(\mu(x))\), and the challenger's belief \(\mu(x)\) at the price \(P(x)\). By
Lemma~\ref{lem:a4-product}, such a tuple is an equilibrium under a pair of tie rules when \((s_+,e(1))\) is a pair of
Lemma~\ref{lem:a4-high} and \(s_-\) is a rate of Lemma~\ref{lem:a4-low} under those rules. Steps 2 to 5 establish these
two conditions.

\emph{Step 2: the quiet market.} Let \(\rho\DT\le2k\), and take \(s_+=s_-=0\) and entry \(\rho\) after every order.
Every posterior equals \(1/2<\tau\). The pair \((0,\rho)\) is the pair (N) of Lemma~\ref{lem:a4-high}, because
\(\rho\DT(1-\pi)=\rho\DT/2\le k\), and the rate \(0\) is the rate (N) of Lemma~\ref{lem:a4-low}, because
\(\rho\DT\pi=\rho\DT/2\le k\). Neither condition involves a tie at the threshold, so the quiet market is an
equilibrium under either challenger rule.

\emph{Step 3: the low-class rate.} Let \(\rho\DT>2k\). By Step 1, \(\rho\DT m\ge k\) holds exactly when
\(b\,(\rho\DT/k-2)\ge\lambda\). At the prior \(1/2\) the rate \(\hat s_-\) of Lemma~\ref{lem:a4-low} is
\(b\,(\rho\DT/k-2)\) when \(\rho\DT m<k<\rho\DT/2\), and \(\lambda\) when \(\rho\DT m\ge k\). Hence
\(\hat s_-=\min\{\lambda,\,b\,(\rho\DT/k-2)\}>0\), the one rate of that lemma, and every tuple in Steps 4 and 5 has
\(s_-=\hat s_-\). Its gross return at entry \(\rho\) satisfies
\begin{equation}\label{eq:o5-sale}
\rho\DT\,\mu_-(\hat s_-)\ge k ,
\end{equation}
with equality when \(\hat s_-<\lambda\). Indeed, if \(\hat s_-=b\,(\rho\DT/k-2)\), then \(2b+\hat s_-=b\rho\DT/k\) and
\(\mu_-(\hat s_-)=k/(\rho\DT)\); if \(\hat s_-=\lambda\), then \(\rho\DT\mu_-(\lambda)=\rho\DT m\ge k\).

\emph{Step 4: a purchase at the rate \(\hat s_-\) does not recruit.} Let \(\rho\DT>2k\) and \(\mu_+(\hat s_-)<\tau\). This
is the case whenever \(\tau\ge1\), because \(\mu_+(\hat s_-)\le M<1\). Take \(s_+=\hat s_-\) and \(e(1)=\rho\). If
\(\hat s_-<\lambda\), then \((\hat s_-,\rho)\) is the pair (Q) of Lemma~\ref{lem:a4-high}: \(0<\hat s_-<\lambda\),
\(\mu_+(\hat s_-)<\tau\), and \(\rho\DT[1-\mu_+(\hat s_-)]=\rho\DT\mu_-(\hat s_-)=k\) by Step 1 and \eqref{eq:o5-sale}. If
\(\hat s_-=\lambda\), then \(M=\mu_+(\lambda)<\tau\), so \((\lambda,\rho)\) is the pair (F), and its return condition
\(\rho\DT(1-M)=\rho\DT m\ge k\) holds. The posteriors \(\mu_-(\hat s_-)<\pi\), \(\mu_0=\pi\) and \(\mu_+(\hat s_-)\) all lie
below \(\tau\), so the tuple is an equilibrium under either challenger rule, with entry \(\rho\) after every order.

\emph{Step 5: a purchase at the rate \(\hat s_-\) recruits.} Let \(\rho\DT>2k\) and \(\mu_+(\hat s_-)\ge\tau\), so that
\(\tau\le\mu_+(\hat s_-)\le M<1\). Take \(e(1)=1\) and \(s_+=\min\{\lambda,\,b\,(\DT/k-2)\}\). By Step 1 and
\eqref{eq:o5-sale},
\begin{equation}\label{eq:o5-buy}
\DT[1-\mu_+(\hat s_-)]=\DT\mu_-(\hat s_-)\ge k/\rho>k .
\end{equation}
Since \(\DT(1-M)=\DT b/(2b+\lambda)\), the inequality \(\DT(1-M)\ge k\) holds exactly when \(b\,(\DT/k-2)\ge\lambda\).

If \(\DT(1-M)\ge k\), then \(s_+=\lambda\ge\hat s_-\), and \(M=\mu_+(\lambda)\ge\mu_+(\hat s_-)\ge\tau\). When \(M>\tau\), the
pair (F) of Lemma~\ref{lem:a4-high} has \(e(1)=1\); when \(M=\tau\), either challenger rule admits \(e(1)=1\). The
return condition \(e(1)\DT(1-M)=\DT(1-M)\ge k\) holds, so \((\lambda,1)\) is the pair (F).

If \(\DT(1-M)<k\), then \(s_+=b\,(\DT/k-2)<\lambda\), and \(\hat s_-<\lambda\), because \(\hat s_-=\lambda\) would give
\(\DT(1-M)=\DT[1-\mu_+(\hat s_-)]>k\) by \eqref{eq:o5-buy}. The return \(\DT[1-\mu_+(s)]=\DT b/(2b+s)\) is continuous
and strictly decreasing in \(s\). It exceeds \(k\) at \(s=\hat s_-\) by \eqref{eq:o5-buy}, and it equals \(k\) exactly
at \(s=s_+\). Hence \(\hat s_-<s_+<\lambda\). Since \(\mu_+\) rises strictly, \(\mu_+(s_+)>\mu_+(\hat s_-)\ge\tau\), so
Lemma~\ref{lem:prices}(d) gives entry one after a purchase under either challenger rule, and \((s_+,1)\) is the pair
(R) of Lemma~\ref{lem:a4-high}.

In both cases \(s_+\ge\hat s_-=s_-\). The posteriors after a sale and after the zero order are at most \(1/2<\tau\),
so their entry is \(\rho\) under either challenger rule, and Step 1 shows that the tuple is an equilibrium.
\end{proof}

\paragraph{Two intervals of coexistence.}
Without the balance condition, the strengths with coexistence along a path can form two disjoint intervals, even
when class \(H\) outbids every incumbent and its value exceeds every incumbent value by more than \(c_H\). Keep every
primitive of the example except the incumbent law, take the prior \(1/2\) and any slot with
\(\reg{tau_r1_frac}<M(1/2)<\reg{tau_r2_frac}\), the example's thresholds at \(r_1\) and \(r_2\), and for
\(r\in[0,2]\) let \(F_r\) put weight \(\reg{twopath_weight_lower}\) on a lower uniform law and the rest on an upper
uniform law, both of half-width \(\reg{twopath_halfwidth}\). The centre of the lower law rises linearly from
\(\reg{twopath_lower_start}\) to \(\reg{twopath_lower_end}\) as \(r\) runs over \([0,1]\) and then stays there; the
centre of the upper law stays at \(\reg{twopath_upper_start}\) on \([0,1]\) and rises linearly to
\(\reg{twopath_upper_end}\) as \(r\) runs over \([1,2]\). Every lower value lies between \(p\) and
\(\ell\), and every upper value between \(\ell\) and \(h\). On \([0,1]\) the incumbent strengthens only where it loses
to both classes: both profits fall, the spread stays constant, \(\tau\) rises and \(J\) falls. On \([1,2]\) it
strengthens between the class values, and the spread widens. A quiet market and an informative market coexist exactly
on \([0,\reg{twopath_lo}]\cup[\reg{twopath_hi},2]\). The strengths with a sure purchase form one interval inside the
second piece, and on the gap the spread lies below \(\DT^{\rm F}\).

% STRUCTURE: the alternative tie rules <- Section 2.4
\paragraph{Other tie rules.}
The results above take the default tie rules. The two alternatives change the region of part~(a) as follows.
\begin{itemize}[leftmargin=2em]
\item When an indifferent high-cost challenger may prepare with any probability, a quiet market and an informative
market coexist against \(F\) exactly when \(F\) satisfies \eqref{eq:region}, or when \(\tau=\pi\) and
\(k/(1-\pi)<\DT\le\DT^{\rm N}\). At \(\tau=\pi\), one quiet market has entry \(\rho\) after every order, and one
informative market is constructed as in part~(a), with entry \(\rho\) after a sale. Inside the region, wherever \(\tau<M\)
and \(J>k\), the threshold pair with \(e(1)=k/J\) is a further informative market, and its entry lies strictly between
\(\rho\) and \eqref{eq:entry-inf}. Where \(\tau=M\), a full purchase is an equilibrium with each entry after a purchase
in \([k/J,1]\).
\item When an indifferent investor abstains and an indifferent high-cost challenger prepares, a quiet market and an
informative market coexist against \(F\) exactly when \(\pi<\tau\le M\) and \(\DT^{\rm F}<\DT\le\DT^{\rm N}\), and the
informative market is the full purchase. When an indifferent investor abstains and an indifferent high-cost challenger
may prepare with any probability, the laws with \(\tau=\pi\) and \(\DT^{\rm F}<\DT\le\DT^{\rm N}\) join this set.
\end{itemize}

\begin{proof}[Proof of the statements on other tie rules]
Each statement follows from Lemmas~\ref{lem:a4-high} to~\ref{lem:a4-product} and the steps of the proof of part~(a).

\emph{An indifferent investor may mix, and an indifferent high-cost challenger may prepare with any probability.} If
\(\tau<\pi\), every posterior of a quiet market lies strictly above \(\tau\), so Step 1 of the proof of part~(a)
applies with no tie, and coexistence fails. If \(\pi<\tau\), Steps 2 and 3 hold as written: the pairs (N) and (Q)
involve no tie at \(\tau\), the pairs that need \(\mu_+(s_+)\ge\tau\) still need \(\tau\le M\) and \(J\ge k\), and the
informative pair constructed in Step 3 remains available, the full purchase with \(e(1)=1\) at \(\tau=M\) and the
threshold pair with \(e^*=1\) at \(J=k\). The coexistence set for \(\pi<\tau\) is therefore \eqref{eq:region}. Inside
the region, where \(\tau<M\) and \(J>k\), the threshold pair has \(0<s^*<\lambda\) and \(e^*=k/J<1\); since \(\tau>\pi\)
and \(\DT\le\DT^{\rm N}\), \(\rho J<\rho\DT(1-\pi)\le k\), so \(e^*>\rho\). With \(s_-=0\), Lemma~\ref{lem:a4-product}
makes it an equilibrium. Its entry \(\rho+(e^*-\rho)(b+\pi s^*)\) exceeds \(\rho\), and it falls short of
\eqref{eq:entry-inf}, which is \(\rho+(1-\rho)(b+\pi s_+)\), because \(e^*<1\) and \(s^*<s_+\), the return at \(s^*\)
being \(J>k\). Where \(\tau=M\), \(J=\DT(1-M)\ge k\) and \(\rho J<k\), and the full purchase meets its return condition
exactly when \(e(1)J\ge k\), that is for \(e(1)\in[k/J,1]\).

Let \(\tau=\pi\). Every quiet market has entry at least \(\rho\), so it needs \(\rho\DT\max\{\pi,1-\pi\}\le k\), that is
\(\DT\le\DT^{\rm N}\); conversely, entry \(\rho\) after every order is allowed at the posterior \(\pi=\tau\), and it gives
a quiet market. Under \(\DT\le\DT^{\rm N}\), a sale at a positive rate has a posterior below \(\tau\), entry \(\rho\) and
a gross return below \(\rho\DT\pi\le k\), so \(s_-=0\) in every equilibrium. A purchase at a positive rate has a
posterior above \(\tau\), so \(e(1)=1\), and its return \(\DT[1-\mu_+(s_+)]\) falls strictly from \(\DT(1-\pi)\) at
\(s_+=0\) to \(\DT(1-M)\) at \(s_+=\lambda\). An informative market therefore exists exactly when \(\DT(1-\pi)>k\): a
full purchase when \(\DT(1-M)\ge k\), and otherwise a mixed purchase at the root in \((0,\lambda)\). Its posteriors
satisfy \(\mu_0<\pi=\mu_-(0)<\mu_+(s_+)\); with entry \(\rho\) after the zero order and, at the tie, after the sale, the
sufficiency argument of Lemma~\ref{lem:a4-product} applies.

\emph{An indifferent investor abstains.} Then (Q), (R) and (T) drop out of Lemma~\ref{lem:a4-high}, the pair (F)
needs a strict return, and the rate (N) of Lemma~\ref{lem:a4-low} is unchanged. Step 1 of the proof of part~(a) uses
strict inequalities only, so it excludes \(\tau<\pi\), and it excludes \(\tau=\pi\) when an indifferent high-cost
challenger prepares. Let \(\pi<\tau\) and \(\DT\le\DT^{\rm N}\). The low class abstains, because a full sale earns
\(\rho\DT m<\rho\DT\pi\le k\), and the informative market is a full purchase with \(\mu_+(\lambda)=M\ge\tau\), entry
one after it, and \(\DT(1-M)>k\), that is \(\DT>\DT^{\rm F}\); a challenger who may prepare with any probability admits
entry one at \(M=\tau\) as well. When it may prepare with any probability and \(\tau=\pi\), the argument for
\(\tau=\pi\) above applies with the full purchase as the only informative pair, so \(\DT^{\rm F}<\DT\le\DT^{\rm N}\).
\end{proof}

\subsubsection*{The example}

\paragraph{The path and its marked strengths.}
In the example the classes are point masses at \(h\) and \(\ell\), the incumbent is uniform on \([0,r]\), and
\((h,\ell,p,\rho,c_H,k)=(\reg{prim_h},\reg{prim_ell},\reg{prim_p},\reg{prim_rho},\reg{prim_cH},\reg{prim_k})\). The path
is \(r\in[r_0,r_2]\), with \((r_0,r_1,r_2)=(\reg{prim_r0},\reg{prim_r1},\reg{prim_r2})\). On it \(\ell<r<h\), so
\(\DT(r)>0\) and \(g_H(r)>g_L(r)\). On the path \(\DT(r)=(r-\ell)^2/(2r)\), with \(\DT'(r)=\tfrac12-\ell^2/(2r^2)>0\), and
the gross return of a purchase that just reaches \(\tau(r)\) and recruits is
\[
J(r)=\DT(r)\,[1-\tau(r)]=\frac{(r-\ell)^2\,[2r(h-c_H)-r^2-p^2]}{2r\,(2rh-r^2-\ell^2)} .
\]
Write \(\fr(y)=\ell+y+\sqrt{y^2+2\ell y}\), the root above \(\ell\) of \(\DT(r)=y\);
\(\DT(r)\le y\) holds exactly when \(r\le\fr(y)\). At the prior \(1/2\) write \(M=M(1/2)\) and \(m=1-M\). Four marked
strengths do not depend on the slot: \(r_A\), the root of \(J(r)=k\) in \((r_0,r_2)\); \(r_N=\fr(2k/\rho)\); \(r_J\),
the root of \(\rho J(r)=k\) in \((r_0,r_2)\); and \(r_*\), the root of \(\rho\DT(r)\tau(r)=k\) in \((r_0,r_2)\). For a
slot with \(0<\lambda<1\), \(0<\beta<1/2\) and \(\tau(r_1)<M<\tau(r_2)\), three more are \(r_F=\fr\bigl(k/(1-M)\bigr)\),
\(r_S=\fr\bigl(k/(\rho m)\bigr)\), and \(r_C\), the root of \(\tau(r)=M\) in \((r_1,r_2)\). Lemma~\ref{lem:a4-path} makes
\(r_A\), \(r_J\), \(r_C\) and \(r_*\) unique.

At \(r_A\) a purchase that just reaches \(\tau\) and recruits earns \(k\); at \(r_J\) the same purchase without
recruitment earns \(k\); at \(r_N\) a unit order against a quiet market earns \(\rho\DT/2=k\); at \(r_F\) a full purchase
that recruits earns \(\DT(1-M)=k\); and at \(r_S\) a full sale earns \(\rho\DT m=k\). The root \(r_C\) of \(\tau(r)=M\),
that is of \(B_r(M)=c_H\), has the closed form
\[
r_C=\frac{Mh-c_H+\sqrt{(Mh-c_H)^2+M[(1-M)\ell^2-p^2]}}{M}.
\]

\begin{lemma}[The example path]\label{lem:a4-path}
On the path \([r_0,r_2]\) of the example:
\begin{enumerate}[label=(\alph*),leftmargin=2em]
\item \(\tau\) is strictly increasing, with \(\tau(r_0)=\reg{tau_r0_frac}>1/2\) and \(\tau(r_2)=\reg{tau_r2_frac}\), so
that \(1-\rho-\tau(r_2)=\reg{margin_tau_r2_3q_B_frac}>0\).
\item The path is a uniform incumbent of Proposition~\ref{prop:coexistence}(c)(iii), with
\(h-r_2-c_H=\reg{margin_unif_afford_B}>0\). The product \(\DT(r)\,[g_H(r)-c_H]\) has derivative at least
\(\reg{margin_balance_B_frac}\cdot\DT'(r)>0\) at every \(r\). Hence \(J\) is strictly increasing, with
\(J(r_0)=\reg{J_r0}<k\) and \(\rho J(r_2)=\reg{rhoJ_r2}>k\).
\item The product \(\DT(r)\,[g_H(r)-c_H]^2\) has derivative at least
\(\reg{margin_balance_sq_B_frac}\cdot[g_H(r)-c_H]\,\DT'(r)>0\) at every \(r\), and \(\DT(r)\,[1-\tau(r)]^2\) is strictly
increasing on the whole path.
\item \(\DT(r)\,\tau(r)\) is strictly increasing.
\end{enumerate}
\end{lemma}

\begin{proof}
On the path \(\ell<r<h\), and Step 1 of the proof of Proposition~\ref{prop:coexistence}(c)(iii) gives
\[
g_H=h-\frac r2-\frac{p^2}{2r},\qquad g_L=\frac{\ell^2-p^2}{2r},\qquad \DT=\frac{(r-\ell)^2}{2r},\qquad
\gamma=\frac{2rh-r^2-\ell^2}{2r}.
\]

(a) The threshold \(\tau=(c_H-g_L)/\gamma\) is the ratio of the numerator \(2rc_H-\ell^2+p^2\) to the denominator
\(2rh-r^2-\ell^2\). The denominator exceeds \(r^2-\ell^2>0\), because \(2h-r>r\). The derivative of \(\tau\) has the sign
of the derivative of the numerator times the denominator, minus the numerator times the derivative of the
denominator, which is
\[
2c_H(r^2-\ell^2)+2(\ell^2-p^2)(h-r)>0 ,
\]
since \(\ell<r<h\) and \(\ell\ge p\). The two values follow from the closed form of \(\tau\), and the last from
\(\rho=\reg{prim_rho}\).

(b) The interval \(I=[r_0,r_2]\) has \(p<\ell<r_0\) and \(h>r_2\), and the margin \(h-r_2-c_H\) is the stated one, so
the path is a uniform incumbent of Proposition~\ref{prop:coexistence}(c)(iii). Directly, \(g_H'=-(r^2-p^2)/(2r^2)\) and
\(\DT/\DT'=r(r-\ell)/(r+\ell)\), so
\[
\bigl(\DT\,(g_H-c_H)\bigr)'=\DT'(g_H-c_H)+\DT g_H'
=\DT'\Bigl[g_H-c_H-\frac{(r^2-p^2)(r-\ell)}{2r(r+\ell)}\Bigr].
\]
In the bracket, \(g_H-c_H\) falls strictly, because \(g_H'<0\) for \(r>p\), and the subtracted term rises strictly, as
the product of the positive increasing functions \((r^2-p^2)/(2r)\) and \((r-\ell)/(r+\ell)\). The bracket is
therefore smallest at \(r_2\), where it equals \(\reg{margin_balance_B_frac}\), and this gives the bound. In
particular \(g_H-c_H>0\) on the path, so the numerator \(\DT(g_H-c_H)\) of \(J=\DT(g_H-c_H)/\gamma\) is positive and
strictly increasing, and \(\gamma\) is positive and strictly decreasing, since \(\gamma'=-\DT'<0\). Hence \(J\) is
strictly increasing. The two values follow from the closed form of \(J\).

(c) In the same way,
\[
\bigl(\DT\,(g_H-c_H)^2\bigr)'=(g_H-c_H)\,\DT'\Bigl[g_H-c_H-\frac{(r^2-p^2)(r-\ell)}{r(r+\ell)}\Bigr],
\]
and the bracket, strictly decreasing for the same reasons, equals \(\reg{margin_balance_sq_B_frac}\) at \(r_2\). Since
\(g_H-c_H>0\), this gives the bound, which is positive. Division by \(\gamma^2\), positive and strictly decreasing,
shows that \(\DT(1-\tau)^2=\DT(g_H-c_H)^2/\gamma^2\) is strictly increasing on the whole path.

(d) \(\DT\) and \(\tau\) are positive and strictly increasing on the path, by \(\DT'>0\) and (a), so their product is
strictly increasing.
\end{proof}

By part (b), \(J(r_0)<k<\rho J(r_2)<J(r_2)\) and \(\rho J(r_0)<k\), so the roots \(r_A\) of \(J=k\) and \(r_J\) of
\(\rho J=k\) exist in \((r_0,r_2)\), and they are unique because \(J\) is strictly increasing. The root \(r_C\) is unique
because \(\tau\) is strictly increasing, and \(r_*\) because \(\DT\tau\) is.

\begin{lemma}[The order of the marked strengths]\label{lem:a4-order}
For every slot with \(0<\lambda<1\), \(0<\beta<1/2\) and \(\tau(r_1)<M(1/2)<\tau(r_2)\),
\[
r_0<r_A<r_F<r_N<r_J<r_S<r_1<r_C<r_2 .
\]
\end{lemma}

\begin{proof}
Each inequality reduces to a fact of Lemma~\ref{lem:a4-path} or to the band \(\tau(r_1)<M<\tau(r_2)\). Every strength
in the chain exceeds \(\ell\), where \(\DT\) is strictly increasing.
\begin{itemize}[leftmargin=2em]
\item \(r_0<r_A<r_1\): \(J(r_0)<k<\rho J(r_1)<J(r_1)\), with \(\rho J(r_1)=\reg{rhoJ_r1}\), and \(J\) is increasing.
\item \(r_A<r_F\): \(\DT(r_A)<k/(1-M)\) is \((1-M)\DT(r_A)<k=(1-\tau(r_A))\DT(r_A)\), that is \(M>\tau(r_A)\). This holds
because \(M>\tau(r_1)>\tau(r_A)\), since \(r_A<r_1\).
\item \(r_F<r_N\): \(k/(1-M)<2k/\rho\) is \(M<1-\rho/2\). This holds because \(M<\tau(r_2)<1-\rho\), by
Lemma~\ref{lem:a4-path}(a).
\item \(r_N<r_J\): \(\rho\DT(r_J)[1-\tau(r_J)]=k\) and \(\tau(r_J)>\tau(r_0)>1/2\) give \(\rho\DT(r_J)>2k\), that is
\(\DT(r_J)>\DT(r_N)\).
\item \(r_J<r_S\): \(\DT(r_J)<k/(\rho m)\) is \(\rho(1-M)\DT(r_J)<k=\rho(1-\tau(r_J))\DT(r_J)\), that is \(M>\tau(r_J)\). The
value \(\rho J(r_1)>k\) gives \(r_J<r_1\), so \(M>\tau(r_1)>\tau(r_J)\).
\item \(r_S<r_1\): the claim is \(\rho m\DT(r_1)>k\). By Lemma~\ref{lem:a4-path}(a), \(m=1-M>1-\tau(r_2)>\rho\), so
\(\rho m\DT(r_1)>\rho^2\DT(r_1)\), and \(\rho^2\DT(r_1)-k=\reg{margin_rho2Delta_r1_B_frac}>0\).
\item \(r_1<r_C<r_2\): this is the band, since \(\tau\) is increasing.\qedhere
\end{itemize}
\end{proof}

The root \(r_*\) lies in \((r_A,r_N)\). Write \(\tau_A=\tau(r_A)\). Since \(J(r_A)=k\), \(\DT(r_A)=k/(1-\tau_A)\), and
\(\rho\DT(r_A)\tau_A<k\) is \(\rho\tau_A<1-\tau_A\), which holds because \((1+\rho)\tau_A<(1+\rho)(1-\rho)<1\) by
Lemma~\ref{lem:a4-path}(a). At \(r_N\), \(\rho\DT(r_N)\tau(r_N)=2k\,\tau(r_N)>k\).

Table~\ref{tab:a4-thresholds} gives the marked strengths at the example's slot, \(\beta=\reg{beta_B}\) and
\(1+u=\eu\), so that \(\lambda=\reg{lambda_B}\). The roots \(r_A\), \(r_J\) and \(r_*\) are unique by
Lemma~\ref{lem:a4-path}; \(r_F\), \(r_N\), \(r_S\) and \(r_C\) have the closed forms above. Each value is rounded to six
places.

\begin{table}[ht]
\centering
\caption{Marked strengths of the example at the prior \(1/2\).}\label{tab:a4-thresholds}
\small
\begin{tabular}{@{}llll@{}}
\toprule
Strength & Definition & Value in the example & Depends on the slot\\
\midrule
\(r_A\) & \(J(r)=k\) & \(\reg{rA}\) & no\\
\(r_F\) & \(\DT(r)(1-M)=k\) & \(\reg{rF_B}\) & yes\\
\(r_N\) & \(\rho\DT(r)/2=k\) & \(\reg{rN}\) & no\\
\(r_J\) & \(\rho J(r)=k\) & \(\reg{rJ}\) & no\\
\(r_S\) & \(\rho\DT(r)m=k\) & \(\reg{rS_B}\) & yes\\
\(r_C\) & \(\tau(r)=M\) & \(\reg{rC_B}\) & yes\\
\(r_*\) & \(\rho\DT(r)\tau(r)=k\) & \(\reg{rstar}\) & no\\
\bottomrule
\end{tabular}
\end{table}

\paragraph{The forms at the prior \(1/2\).}
Take the prior \(1/2\) and a slot with \(\tau(r_1)<M<\tau(r_2)\). Then \(1-\mu_+(s_+)=b/(2b+s_+)\),
\(\mu_-(s_-)=b/(2b+s_-)\) and \(s^*=b(2\tau-1)/(1-\tau)\). Under the default tie rules the low class sells at the rate
of Lemma~\ref{lem:a4-low},
\[
\hat s_-(r)=0\ \text{on}\ [r_0,r_N],\qquad \hat s_-(r)=b\,[\rho\DT(r)/k-2]\ \text{on}\ (r_N,r_S),\qquad
\hat s_-(r)=\lambda\ \text{on}\ [r_S,r_2],
\]
and class \(H\)'s pair takes one of five forms:
\begin{itemize}[leftmargin=2em]
\item \textbf{N}, a quiet market: \(s_+=0\) and \(e(1)=\rho\);
\item \textbf{R}, a mixed purchase that recruits: \(s_+=s_R(r):=b[\DT(r)/k-2]\) and \(e(1)=1\);
\item \textbf{F}, a full purchase: \(s_+=\lambda\), with \(e(1)=1\) for \(r<r_C\), \(e(1)=1\) at \(r_C\) because an
indifferent high-cost challenger prepares, and \(e(1)=\rho\) for \(r>r_C\);
\item \textbf{Q}, mixed orders at the entry floor: \(s_+=s_Q(r):=b[\rho\DT(r)/k-2]=\hat s_-(r)\) and \(e(1)=\rho\);
\item \textbf{T}, the threshold pair: \(s_+=s^*(r)\) and \(e(1)=e^*(r)=k/J(r)\), which requires \(J(r)=k\), that is
\(r=r_A\).
\end{itemize}
The forms R, Q and T need an indifferent investor who may mix. When an indifferent high-cost challenger may prepare
with any probability, a full purchase at \(r_C\) admits every \(e(1)\in[\rho,1]\), and the threshold pair is a pair of
Lemma~\ref{lem:a4-high} exactly on \([r_A,r_J]\), where \(\rho J(r)\le k\le J(r)\).

Each form defines the equilibrium tuple completely:
\begin{itemize}[leftmargin=2em]
\item orders: \(\sigma_H(1)=s_+/\lambda\), \(\sigma_H(0)=1-s_+/\lambda\), \(\sigma_L(-1)=\hat s_-(r)/\lambda\),
\(\sigma_L(0)=1-\hat s_-(r)/\lambda\);
\item posteriors: \(\mu_-(\hat s_-)\), \(\mu_0(s_+,\hat s_-)\) and \(\mu_+(s_+)\) after the orders \(-1\), \(0\) and \(1\);
\item challenger beliefs: \(\bar\mu(P(x))=\mu(x)\) at each price;
\item high-cost preparation: \(\zeta(P(1))=(e(1)-\rho)/(1-\rho)\), and \(\zeta(P(x))=0\) at the prices of the
orders \(0\) and \(-1\);
\item prices: \(P(x)=t_0+e(x)\,w(\mu(x))\);
\item bids: truthful.
\end{itemize}
If two orders share a price, they share the posterior and the entry (Lemma~\ref{lem:a4-product}), so \(\zeta\) is
well defined as a function of the price.

% STRUCTURE: full statement of Corollary 2 <- Section 4.3
\begin{named}{Corollary 2 (Coexistence in the example, full statement)}\fixedlabel{cor:example-coexistence}{2}
Take the example at the prior \(1/2\), a slot with \(0<\lambda<1\), \(0<\beta<1/2\) and \(\tau(r_1)<M(1/2)<\tau(r_2)\),
and the default tie rules. The equilibrium set at each strength of the path is
\begin{center}
\begin{tabular}{@{}ll@{}}
\toprule
Strength & Equilibria\\
\midrule
\(r_0\le r<r_A\) & \textbf{N}\\
\(r=r_A\) & \textbf{N}; \textbf{T} with entry one after a purchase\\
\(r_A<r<r_F\) & \textbf{N}; \textbf{R}\\
\(r_F\le r\le r_N\) & \textbf{N}; \textbf{F}\\
\(r_N<r<r_J\) & \textbf{F}; \textbf{Q}\\
\(r_J\le r\le r_2\) & \textbf{F}\\
\bottomrule
\end{tabular}
\end{center}
In each, the low class sells at the rate \(\hat s_-(r)\). Entry is \(\rho\) in a quiet market and in mixed orders at the
entry floor; \(\Ent_{\rm R}(r)=\rho+(1-\rho)\,b\DT(r)/(2k)\) in a mixed purchase that recruits, and \(\Ent_{\rm R}(r_A)\)
in the threshold pair at \(r_A\); and \(\Ent_{\rm F}=\rho+(1-\rho)(b+\lambda/2)\) in a full purchase up to \(r_C\), with
\(\Ent_{\rm F}=\rho\) above \(r_C\).
\begin{enumerate}[label=(\alph*),leftmargin=2em]
\item A quiet market and an informative market coexist exactly on \([r_A,r_N]\).
\item Along the informative branch, entry is \(\rho+(1-\rho)\min\{b\DT(r)/(2k),\,b+\lambda/2\}\) on \([r_A,r_C]\) and
\(\rho\) on \((r_C,r_2]\). It is continuous on \([r_A,r_C]\), rises strictly on \([r_A,r_F]\), is constant on
\([r_F,r_C]\), and returns to the floor \(\rho\) above the entry ceiling \(r_C\).
\item The equilibrium set is nonempty at every strength.
\end{enumerate}
\end{named}

\begin{proof}[Proof of Corollary~\ref{cor:example-coexistence}]
Fix a slot in the band and the prior \(1/2\). The proof has eight steps. Steps 2 and 3 classify the pairs and the
rates under each pair of tie rules; Step 5 assembles the set under the default tie rules.

\emph{Step 1: the frame.} By Lemma~\ref{lem:a4-path}(a), \(\tau(r)\ge\tau(r_0)>1/2=\pi\) on the path, so
Lemmas~\ref{lem:a4-high} to~\ref{lem:a4-product} apply at every \(r\). Every equilibrium has a high-class pair
from Lemma~\ref{lem:a4-high} and a low-class rate from Lemma~\ref{lem:a4-low}. Step 4 shows \(s_+\ge s_-\) for every
such combination, and Lemma~\ref{lem:a4-product} then gives the equilibrium set as their product. Recall
\(1-M=1/(2+u)\) and \(ub=\lambda\).

\emph{Step 2: the high-class pairs.} Use Lemma~\ref{lem:a4-path}, the definitions of the marked strengths, and
Lemma~\ref{lem:a4-order}.
\begin{itemize}[leftmargin=2em]
\item (N) holds if and only if \(\rho\DT/2\le k\), that is \(r\le r_N\).
\item (R): \(\DT b/(2b+s_+)=k\) gives \(s_+=s_R(r)\). Then \(\mu_+(s_R)>\tau\) if and only if \(k/\DT<1-\tau\), that is
\(J(r)>k\), that is \(r>r_A\). Also \(s_R<\lambda\) if and only if \(\DT/k<2+u\), that is \(\DT(1-M)<k\), that is
\(r<r_F\). So (R) holds if and only if \(r_A<r<r_F\), with an indifferent investor who may mix.
\item (Q): \(\rho\DT b/(2b+s_+)=k\) gives \(s_+=s_Q(r)\). Then \(s_Q>0\) if and only if \(r>r_N\); \(\mu_+(s_Q)<\tau\) if
and only if \(\rho J(r)<k\), that is \(r<r_J\); and \(s_Q<\lambda\) if and only if \(\rho\DT(1-M)<k\), that is
\(r<r_S\). Since \(r_J<r_S\), (Q) holds if and only if \(r_N<r<r_J\), with an indifferent investor who may mix.
\item (T): \(s^*\in(0,\lambda)\) if and only if \(1/2<\tau(r)<M\), that is \(r<r_C\). Also \(\rho\le e^*\le1\) if and
only if \(\rho J(r)\le k\le J(r)\), that is \(r_A\le r\le r_J\), and \(e^*=1\) if and only if \(r=r_A\). Since
\(r_J<r_C\), (T) holds at \(r_A\) when an indifferent high-cost challenger prepares, and on \([r_A,r_J]\) when it may
prepare with any probability, with an indifferent investor who may mix in both cases.
\item (F): for \(r<r_C\), \(M>\tau(r)\), so \(e(1)=1\), and the condition is \(\DT(1-M)\ge k\), that is \(r\ge r_F\)
(\(r>r_F\) when an indifferent investor abstains). For \(r>r_C\), \(M<\tau(r)\), so \(e(1)=\rho\), and
\(\rho\DT(1-M)>k\) holds because \(r>r_C>r_S\). At \(r=r_C\), an indifferent high-cost challenger who prepares gives
\(e(1)=1\); one who may prepare with any probability allows any \(e(1)\in[\rho,1]\), and every such value meets the
condition because \(\rho\DT(r_C)(1-M)>k\). So (F) holds if and only if \(r\ge r_F\) when an indifferent investor may
mix, and \(r>r_F\) when it abstains.
\end{itemize}

\emph{Step 3: the low-class rate.} At the prior \(1/2\), Lemma~\ref{lem:a4-low} gives, when an indifferent investor
may mix, \(s_-=0\) if and only if \(r\le r_N\), \(s_-=\lambda\) if and only if \(r\ge r_S\), and
\(\hat s_-(r)=b(\rho\DT/k-2)\) on \((r_N,r_S)\). When an indifferent investor abstains, no rate qualifies if and only
if \(\rho\DT m\le k<\rho\DT/2\), that is on \((r_N,r_S]\).

\emph{Step 4: \(s_+\ge s_-\).} For \(r\le r_N\), \(s_-=0\le s_+\). For \(r>r_N\), the pairs are (F) with
\(s_+=\lambda\ge s_-\); (Q) with \(s_+=s_Q(r)=\hat s_-(r)\); and (T) with \(r_N<r\le r_J\). For (T),
\(s^*\ge\hat s_-(r)\) if and only if \((2\tau-1)/(1-\tau)\ge\rho\DT/k-2\), that is \(1/(1-\tau)\ge\rho\DT/k\), that is
\(\rho J(r)\le k\), which holds for \(r\le r_J\). By Lemma~\ref{lem:a4-product}, entry after the zero order and after
a sale is \(\rho\) in every equilibrium.

\emph{Step 5: the set.} Combine Steps 2 to 4 with the order of Lemma~\ref{lem:a4-order}. Under the default tie rules
the set is: a quiet market on \([r_0,r_A)\); a quiet market and the threshold pair with \(e^*=1\) at \(r_A\); a quiet
market and a mixed purchase that recruits on \((r_A,r_F)\); a quiet market and a full purchase on \([r_F,r_N]\); a full
purchase and mixed orders at the entry floor on \((r_N,r_J)\); and a full purchase on \([r_J,r_2]\). In each, the low
class sells at the rate \(\hat s_-(r)\) of Step 3. On \([r_0,r_N]\) this is Proposition~\ref{prop:coexistence}(a) on the
path, since there \(\DT(r)\le\DT^{\rm N}\), \(\pi<\tau(r)<M\), and \(J(r)\ge k\) exactly when \(r\ge r_A\).

\emph{Step 6: the tuples.} Fix a form and a strength at which Steps 2 and 3 admit it. Step 4 gives \(s_+\ge s_-\), so
the sufficiency part of Lemma~\ref{lem:a4-product} shows that the tuple satisfies (E1) to (E6). In particular, if two
orders have equal posteriors, one of them is \(0\) or \(-1\), so the common posterior is at most \(1/2<\tau\), and both
orders have entry \(\rho\) and the same price.

\emph{Step 7: entry.} Entry is \(\Ent=\sum_x\Pr(X=x)\,e(x)\). Since \(e(0)=e(-1)=\rho\) and
\(\Pr(X=1)=\tfrac12(b+s_+)+\tfrac12b=b+s_+/2\), \(\Ent=\rho+[e(1)-\rho](b+s_+/2)\).
\begin{itemize}[leftmargin=2em]
\item A mixed purchase that recruits: \(b+s_R/2=b\DT/(2k)\) and \(e(1)=1\), so
\(\Ent_{\rm R}(r)=\rho+(1-\rho)\,b\DT(r)/(2k)\). It rises strictly in \(r\) because \(\DT\) does. At \(r_F\),
\(\DT(r_F)=k/(1-M)=k(2+u)\), so \(b\DT(r_F)/(2k)=b+ub/2=b+\lambda/2\), and the mixed purchase joins the full purchase
continuously.
\item A full purchase: \(b+\lambda/2\), with \(e(1)=1\) up to \(r_C\) and \(e(1)=\rho\) above \(r_C\). So
\(\Ent_{\rm F}=\rho+(1-\rho)(b+\lambda/2)\) on \([r_F,r_C]\) and \(\Ent_{\rm F}=\rho\) on \((r_C,r_2]\).
\item A quiet market and mixed orders at the entry floor: \(e(1)=\rho\), so \(\Ent=\rho\).
\item The threshold pair at \(r_A\): \(J(r_A)=k\) gives \(\DT(r_A)/k=1/(1-\tau(r_A))\), so \(s^*(r_A)=s_R(r_A)\), and with
\(e^*=1\) its entry is \(\Ent_{\rm R}(r_A)\).
\end{itemize}
Along the informative branch, entry is therefore \(\rho+(1-\rho)\min\{b\DT(r)/(2k),\,b+\lambda/2\}\) on \([r_A,r_C]\): it is
continuous, rises strictly on \([r_A,r_F]\) because \(\DT\) does, and is constant on \([r_F,r_C]\). Above \(r_C\) it is
\(\rho\). The strict rise uses a mixed purchase, which needs an indifferent investor who may mix.

\emph{Step 8: existence.} Proposition~\ref{prop:coexistence}(d) applies at every strength, since \(\tau(r)>1/2\) by
Step 1. Directly, a quiet market exists on \([r_0,r_N]\) and a full purchase on \([r_F,r_2]\), and \(r_F<r_N\).

Parts (a), (b) and (c) are Steps 5, 7 and 8.
\end{proof}

\paragraph{The example's slot.}
Table~\ref{tab:a4-set} gives the equilibrium set at the example's slot under the default tie rules. In a mixed purchase
that recruits, entry is \(\reg{entry_mb_r140_B}\) at \(r=\reg{str_r140}\) and \(\reg{entry_mb_r145_B}\) at
\(r=\reg{str_r145}\). The low class abstains on \([r_0,r_N]\), mixes on \((r_N,r_S)\) and sells for sure on
\([r_S,r_2]\). The strengths \(\reg{str_r155}\), \(\reg{str_r160}\) and \(\reg{str_r165}\) lie in \([r_F,r_N]\), where
the equilibria are a quiet market and a full purchase under the default tie rules. When an indifferent high-cost
challenger may prepare with any probability, the threshold pair adds entry \(\reg{entry_th_r155_B}\),
\(\reg{entry_th_r160_B}\) and \(\reg{entry_th_r165_B}\) at these three strengths. The first two values are
\(\Ent_{\rm R}\) of Corollary~\ref{cor:example-coexistence}, and the last three are \(\Ent_{\rm T}\) of
Proposition~\ref{prop:coexistence}(c)(iv) at the prior \(1/2\), each at the stated strength.

\begin{table}[ht]
\centering
\caption{The equilibrium set of the example at the prior \(1/2\) and the example's slot, under the default tie
rules.}\label{tab:a4-set}
\small
\begin{tabular}{@{}lll@{}}
\toprule
Strengths & Equilibria & Entry\\
\midrule
\([r_0,\,\reg{rA})\) & \textbf{N} & \(\reg{entry_r0}\)\\
\(\reg{rA}\) & \textbf{N}; \textbf{T} with \(e^*=1\) & \(\reg{entry_r0}\); \(\reg{entry_rA_B}\)\\
\((\reg{rA},\,\reg{rF_B})\) & \textbf{N}; \textbf{R} & \(\reg{entry_r0}\); from \(\reg{entry_rA_B}\) up to \(\reg{entry_r1_B}\)\\
\([\reg{rF_B},\,\reg{rN}]\) & \textbf{N}; \textbf{F} & \(\reg{entry_r0}\); \(\reg{entry_r1_B}\)\\
\((\reg{rN},\,\reg{rJ})\) & \textbf{F}; \textbf{Q} & \(\reg{entry_r1_B}\); \(\reg{entry_r0}\)\\
\([\reg{rJ},\,\reg{rC_B}]\) & \textbf{F} & \(\reg{entry_r1_B}\)\\
\((\reg{rC_B},\,r_2]\) & \textbf{F} & \(\reg{entry_r2}\)\\
\bottomrule
\end{tabular}
\end{table}

% STRUCTURE: Corollary A.2 and its definitions <- Section 4.3
\paragraph{Coexistence at other priors.}
For a prior \(\pi\), write \(r_N(\pi)=\fr\bigl(k/(\rho\max\{\pi,1-\pi\})\bigr)\). When some strength on the path has
\(\tau(r)\le M(\pi)\), let \(r_C(\pi)\) be the largest one. When \(\tau(r_0)\le\pi\le\tau(r_2)\), let \(r_\pi\) solve
\(\tau(r_\pi)=\pi\). Write \(\tau_A=\tau(r_A)\). Three priors divide the cases: \(\pi_A^{\rm lo}\) solves
\(M(\pi_A^{\rm lo})=\tau_A\), that is \(\pi_A^{\rm lo}=\tau_A/[\tau_A+(1-\tau_A)(1+u)]\); \(\pi_A^{\rm hi}=\tau_A\); and
\(\pi_{\max}=\tau(r_*)\).

\begin{corollary}[Coexistence at every prior]\label{prop:priors}
Take the example, any slot with \(0<\lambda<1\) and \(0<\beta<1/2\), any prior \(\pi\in(0,1)\), and the default tie
rules. The set \(\mathcal C(\pi)\) of strengths on the path at which a quiet market and an informative market coexist
is the set of Proposition~\ref{prop:coexistence}(a) on the path,
\[
\mathcal C(\pi)=\bigl\{r\in[r_0,r_2]:\ r\ge r_A,\ \ \pi<\tau(r)\le M(\pi),\ \ \rho\DT(r)\max\{\pi,1-\pi\}\le k\bigr\}.
\]
\begin{enumerate}[label=(\alph*),leftmargin=2em]
\item If \(\pi<\pi_A^{\rm lo}\), \(\mathcal C(\pi)\) is empty. At \(\pi=\pi_A^{\rm lo}\), \(\mathcal C(\pi)=\{r_A\}\).
\item If \(\pi_A^{\rm lo}<\pi<\pi_A^{\rm hi}\), then \(\mathcal C(\pi)=[r_A,\min\{r_N(\pi),r_C(\pi)\}]\), a nondegenerate
interval. Coexistence starts at \(r_A\) for every prior in this window and every slot.
\item If \(\pi_A^{\rm hi}\le\pi<\pi_{\max}\), then \(\mathcal C(\pi)=(r_\pi,\min\{r_N(\pi),r_C(\pi)\}]\), a nonempty
interval with \(r_\pi\ge r_A\). Coexistence starts where the prior reaches the threshold.
\item If \(\pi\ge\pi_{\max}\), \(\mathcal C(\pi)\) is empty.
\end{enumerate}
At each \(r\in\mathcal C(\pi)\) the informative market is a mixed purchase that recruits or a full purchase, and the
threshold pair at \(r_A\) in case (b). Its entry is \(\rho+(1-\rho)\min\{(1-\pi)\,b\DT(r)/k,\ b+\pi\lambda\}\), which
rises strictly while the purchase is mixed and is then constant.
\end{corollary}

\begin{proof}[Proof of Corollary~\ref{prop:priors}]
Fix a slot and a prior. Every law of the path satisfies the hypotheses of Proposition~\ref{prop:coexistence}, since
\(\ell<r<h\). By Proposition~\ref{prop:coexistence}(a), coexistence holds at \(r\) if and only if
\(\pi<\tau(r)\le M(\pi)\) and \(k/(1-\tau(r))\le\DT(r)\le\DT^{\rm N}\). Under \(\tau(r)\le M(\pi)<1\), the first spread
inequality is \(J(r)\ge k\), which holds exactly when \(r\ge r_A\) by Lemma~\ref{lem:a4-path}(b), and the second is
\(\rho\DT(r)\max\{\pi,1-\pi\}\le k\). This gives the display. By the same part, the informative market is a full
purchase where \(\DT(r)[1-M(\pi)]\ge k\), the threshold pair with entry one at \(r_A\) where
\(\DT(r_A)[1-M(\pi)]<k\), and a mixed purchase that recruits elsewhere; its entry is \eqref{eq:entry-inf} at
\(\DT=\DT(r)\), which rises strictly while the purchase is mixed, because \(\DT\) rises strictly on the path, and is then
constant.

By Lemma~\ref{lem:a4-path}(a), \(\pi<\tau(r)\) holds for every \(r\) if \(\pi<\tau(r_0)\), for \(r>r_\pi\) if
\(\tau(r_0)\le\pi\le\tau(r_2)\), and for no \(r\) if \(\pi>\tau(r_2)\). The condition \(\tau(r)\le M(\pi)\) holds for
\(r\le r_C(\pi)\) when \(r_C(\pi)\) exists, and the last condition holds for \(r\le r_N(\pi)\). Also
\(\rho\DT(r_A)\max\{\pi,1-\pi\}<\rho\DT(r_A)=\rho k/(1-\tau_A)<k\), because \(\tau_A<\tau(r_2)<1-\rho\) by
Lemma~\ref{lem:a4-path}(a). So \(r_A<r_N(\pi)\) for every prior.
\begin{itemize}[leftmargin=2.5em]
\item[(a)] \(M\) is increasing in \(\pi\), so \(\pi<\pi_A^{\rm lo}\) gives \(M(\pi)<\tau_A\le\tau(r)\) for every
\(r\ge r_A\), and no strength qualifies. At \(\pi=\pi_A^{\rm lo}<\tau_A\) only \(r_A\) qualifies, and there
\(M(\pi)=\tau_A\): the informative market is the full purchase, with \(\DT(r_A)(1-\tau_A)=k\).
\item[(b)] Let \(\pi_A^{\rm lo}<\pi<\pi_A^{\rm hi}\); the window is nonempty because \(u>0\). For \(r\ge r_A\),
\(\tau(r)\ge\tau_A>\pi\). Also \(\tau_A<M(\pi)\) gives \(r_A<r_C(\pi)\), and \(r_A<r_N(\pi)\). So
\(\mathcal C(\pi)=[r_A,\min\{r_N(\pi),r_C(\pi)\}]\), a nondegenerate interval. At \(r_A\),
\(\DT(r_A)[1-M(\pi)]<\DT(r_A)(1-\tau_A)=k\), so the informative market is the threshold pair with entry one. The start
\(r_A\) does not depend on the prior in this window or on the slot.
\item[(c)] Let \(\pi_A^{\rm hi}\le\pi<\pi_{\max}\). The window is nonempty, because \(r_A<r_*\) and \(\tau\) rises
strictly. Here \(\tau(r_0)<\tau_A\le\pi<\tau(r_*)<\tau(r_2)\), so \(r_\pi\) exists, \(r_\pi\ge r_A\), and
\(\pi<\tau(r)\) holds exactly for \(r>r_\pi\). The interval \(\mathcal C(\pi)=(r_\pi,\min\{r_N(\pi),r_C(\pi)\}]\) is
nonempty: \(r_\pi<r_C(\pi)\) because \(\tau(r_\pi)=\pi<M(\pi)\); and since \(\pi>1/2\), \(r_\pi<r_N(\pi)\) if and only if
\(\rho\DT(r_\pi)\tau(r_\pi)<k\), that is \(r_\pi<r_*\) by Lemma~\ref{lem:a4-path}(d), that is \(\pi<\pi_{\max}\).
\item[(d)] Let \(\pi\ge\pi_{\max}\). Then \(\pi>1/2\), and any \(r\) with \(\pi<\tau(r)\) has \(r>r_\pi\ge r_*\), so
\(\rho\DT(r)\pi>\rho\DT(r_*)\tau(r_*)=k\), and no strength qualifies.\qedhere
\end{itemize}
\end{proof}

The priors \(\pi_A^{\rm hi}\) and \(\pi_{\max}\) do not depend on the slot. In the example
\(\pi_A^{\rm hi}\approx\reg{piA_hi}\) and \(\pi_{\max}\approx\reg{pimax}\), and at the example's slot
\(\pi_A^{\rm lo}\approx\reg{piA_lo_B}\).

\paragraph{Other tie rules in the example.}
Take the prior \(1/2\) and a slot with \(\tau(r_1)<M(1/2)<\tau(r_2)\). The threshold pair has entry one after a
purchase at \(r_A\) and is an equilibrium there under the default tie rules. When an indifferent investor may mix and
an indifferent high-cost challenger may prepare with any probability, the threshold pair is also an equilibrium on
\((r_A,r_J]\), and its entry falls strictly from \(\Ent_{\rm R}(r_A)\) at \(r_A\) to \(\rho\) at \(r_J\), by
Lemma~\ref{lem:a4-path}(b) and~(c) and Proposition~\ref{prop:coexistence}(c)(iv); at \(r_C\), entry after a purchase in
a full purchase can then take any value in \([\rho,1]\). When an indifferent investor abstains, under either
challenger rule, the equilibrium set is a quiet market on \([r_0,r_F]\); a quiet market and a full purchase on
\((r_F,r_N]\); empty on \((r_N,r_S]\), where the low class has no qualifying sale rate; and a full purchase on
\((r_S,r_2]\).

At a general prior and any slot with \(0<\lambda<1\) and \(0<\beta<1/2\), when an indifferent investor may mix and an
indifferent high-cost challenger may prepare with any probability, the coexistence set is
\[
\widetilde{\mathcal C}(\pi)=\mathcal C(\pi)\cup\{r_\pi\}\quad\text{for }\tau_A<\pi\le\pi_{\max},\qquad
\widetilde{\mathcal C}(\pi)=\mathcal C(\pi)\ \text{otherwise}.
\]
At \(\pi=\pi_{\max}\), the point \(r_\pi=r_*=\reg{rstar}\) is the whole set. At the example's slot
(\(\lambda=\reg{lambda_B}\), \(\beta=\reg{beta_B}\)), the top posterior is \(M(\pi_{\max})=\reg{M_pimax_B}\), the
posterior after the zero order is \(\reg{mu0_pimax_B}\), the informative market is a full purchase, which earns
\(\DT(r_*)[1-M(\pi_{\max})]-k=\reg{margin_cmend_B}\) net of the trading cost, and entry is \(\reg{entry_cmend_B}\). At every
prior and every slot, when an indifferent investor abstains and an indifferent high-cost challenger prepares,
coexistence holds on the strengths of \(\mathcal C(\pi)\) above \(\fr\bigl(k/(1-M(\pi))\bigr)\); when an indifferent
investor abstains and an indifferent high-cost challenger may prepare with any probability, the point \(r_\pi\) joins
them if \(\fr\bigl(k/(1-M(\pi))\bigr)<r_\pi\le r_N(\pi)\). Table~\ref{tab:a4-ties} records the dependence at the
prior \(1/2\).

\begin{proof}[Proof of the statements on other tie rules in the example]
At the prior \(1/2\), Steps 2 to 4 of the proof of Corollary~\ref{cor:example-coexistence} hold under each pair of tie
rules. Under the default tie rules, Step 2 gives the threshold pair at \(r_A\) with \(e^*=1\), and Step 7 gives its
entry \(\Ent_{\rm R}(r_A)\). When an indifferent investor may mix and an indifferent high-cost challenger may prepare
with any probability, Step 2 gives the threshold pair on \([r_A,r_J]\), Step 4 gives \(s_+\ge s_-\), and
Lemma~\ref{lem:a4-product} completes the tuple. On the path the strengths of Proposition~\ref{prop:coexistence}(c)(iv)
are \([r_A,r_J]\): \(\tau>1/2\) everywhere, \(\tau<M\) exactly below \(r_C\), \(r_J<r_C\), and \(\rho J\le k\le J\) exactly
on \([r_A,r_J]\). Lemma~\ref{lem:a4-path}(b) and~(c) give the hypotheses of that part, and they give strict changes as
well: \(k-\rho J\) is positive on \([r_A,r_J)\), zero at \(r_J\) and strictly decreasing, and \(\DT(1-\tau)^2\) is positive
and strictly increasing, so \(\Ent_{\rm T}\) falls strictly. At \(r_A\), \(e^*=1\) gives \(\Ent_{\rm T}(r_A)=\Ent_{\rm R}(r_A)\);
at \(r_J\), \(\rho J=k\) gives \(\Ent_{\rm T}(r_J)=\rho\). The full purchase at \(r_C\) is in Step 2. When an indifferent
investor abstains, Step 2 removes the pairs that need a mixing investor and leaves the full purchase on
\((r_F,r_2]\), and Step 3 leaves no low-class rate on \((r_N,r_S]\); the stated set follows.

At a general prior, the statements on other tie rules against a general law apply at each strength of the path. When
an indifferent investor may mix and an indifferent high-cost challenger may prepare with any probability, they add to
\(\mathcal C(\pi)\) the strengths with \(\tau(r)=\pi\) and \(k/(1-\pi)<\DT(r)\le\DT^{\rm N}\). Such a strength is \(r_\pi\).
There \(k/(1-\pi)<\DT(r_\pi)\) is \(J(r_\pi)>k\), that is \(r_\pi>r_A\), or \(\pi>\tau_A\), by Lemma~\ref{lem:a4-path}(a)
and~(b); and for \(\pi>\tau_A>1/2\), \(\DT(r_\pi)\le\DT^{\rm N}\) is \(\rho\DT(r_\pi)\tau(r_\pi)\le k\), that is
\(r_\pi\le r_*\), or \(\pi\le\pi_{\max}\), by Lemma~\ref{lem:a4-path}(d). For \(\tau_A<\pi\le\pi_{\max}\) the root
\(r_\pi\) exists, because \(\tau(r_0)<\tau_A\) and \(\tau(r_*)<\tau(r_2)\). At \(\pi_{\max}\), Corollary~\ref{prop:priors}(d)
makes \(\mathcal C(\pi_{\max})\) empty, so the point \(r_*\) is the whole set. At the example's slot the four values
follow from the construction at \(\tau=\pi\), with the purchase rate, the posteriors and the entry formula of
Proposition~\ref{prop:coexistence}(a); the purchase is full because \(\DT(r_*)[1-M(\pi_{\max})]>k\). When an
indifferent investor abstains, the general statements give the strengths of \(\mathcal C(\pi)\) with
\(\DT(r)>k/(1-M(\pi))\), that is \(r>\fr\bigl(k/(1-M(\pi))\bigr)\), when the challenger prepares; when it may prepare
with any probability, \(r_\pi\) joins them if \(k/(1-M(\pi))<\DT(r_\pi)\le\DT^{\rm N}\), that is
\(\fr\bigl(k/(1-M(\pi))\bigr)<r_\pi\le r_N(\pi)\).
\end{proof}

\begin{table}[ht]
\centering
\caption{Which conclusions depend on the tie rules, at the prior \(1/2\). The column heads give the choice of an
indifferent investor, which may mix over its best orders or abstains, and of an indifferent high-cost challenger,
which prepares or may prepare with any probability.}\label{tab:a4-ties}
\small
\begin{tabular}{@{}lllll@{}}
\toprule
 & \multicolumn{2}{l}{Investor may mix} & \multicolumn{2}{l}{Investor abstains}\\
\cmidrule(lr){2-3}\cmidrule(lr){4-5}
Result & prepares & any probability & prepares & any probability\\
\midrule
Proposition~\ref{prop:ccc} & holds & holds & holds & holds\\
Quiet market \textbf{N} & \([r_0,r_N]\) & \([r_0,r_N]\) & \([r_0,r_N]\) & \([r_0,r_N]\)\\
Mixed purchase \textbf{R} & \((r_A,r_F)\) & \((r_A,r_F)\) & none & none\\
Threshold pair \textbf{T} & \(r_A\) only & \([r_A,r_J]\) & none & none\\
Full purchase \textbf{F} & \([r_F,r_2]\) & \([r_F,r_2]\) & \((r_F,r_N]\cup(r_S,r_2]\) & \((r_F,r_N]\cup(r_S,r_2]\)\\
Set on \((r_N,r_J)\) & \textbf{F}; \textbf{Q} & \textbf{F}; \textbf{Q}; \textbf{T} & empty & empty\\
Set at \(r_J\) & \textbf{F} & \textbf{F}; \textbf{T} & empty & empty\\
Set on \((r_J,r_S]\) & \textbf{F} & \textbf{F} & empty & empty\\
Entry after a purchase at \(r_C\) & \(1\) & any in \([\rho,1]\) & \(1\) & any in \([\rho,1]\)\\
Coexistence & \([r_A,r_N]\) & \([r_A,r_N]\) & \((r_F,r_N]\) & \((r_F,r_N]\)\\
\bottomrule
\end{tabular}
\end{table}

When an indifferent investor abstains, a full purchase on \((r_N,r_S]\) is not an equilibrium, because the low class
has no qualifying rate there (Lemma~\ref{lem:a4-low}). So the equilibrium set is nonempty along the whole path when an
indifferent investor may mix, and not when it abstains; the strict rise of entry along a mixed purchase that recruits
also uses the first rule. The challenger rule matters only through the threshold pair and the point \(r_C\).
